Radio frequency identification (RFID) has been widely studied as part of the broader class of automatic identification technologies. Kaur et al. provides a comprehensive overview of its principles, capabilities, and limitations. RFID systems enable the wireless transmission of identification using radio waves, allowing individuals or objects to be identified without requiring line-of-sight interaction, in contrast to traditional barcode systems [1]. This capability significantly improves automation and reduces the need for human action within lock systems, making RFID useful for embedded and access control applications that require high speed and reliability. 

A major focus of prior work is the distinction between passive and active RFID systems, as well as the physical mechanisms that enable their operation. Passive RFID operates by harvesting energy from the electromagnetic field generated by a reader. In near-field systems, this process is governed by magnetic induction, where an alternating magnetic field induces a voltage across a tag's coil, which provides enough energy for the tag to transmit its stored ID [1]. This eliminates the need for an onboard power source. These characteristics make passive RFID well-suited for compact, power-constrained embedded systems, such as our intended implementation within a firearm's firing control group. 

Despite these advantages, Kaur et al. also identifies several important limitations of RFID systems that must be addressed. One key constraint is the relationship between operating frequency and read range. In near-field systems, the effective communication range is limited and decreases rapidly due to the inverse cubic decay of the magnetic field [1]. While this restricts long-range communication, it can be advantageous for security sensitive applications by ensuring that the authenticated user is within close proximity. Other challenges include signal collision when multiple tags are present, environmental interference caused by conductive materials or electromagnetic noise, and the lack of universal standardization across RFID systems [1]. Security and privacy concerns are also significant, as unauthorized cloning of tags can compromise system integrity if appropriate safeguards are not implemented.

Taken together, the existing literature establishes RFID as a mature and well-understood technology for identification and access control. It has demonstrated advantages in automation, robustness, and low power operation. At the same time, it highlights key technical and security challenges that must be carefully addressed in our implementation. This paper provides a strong foundation for the proposed system, informing both the design choices and the tradeoffs involved in applying RFID to firearm safety.
